% Chinese production target: zh-Hans-CN mathematical register.
% Source authority: sealed German P31 cumulative TeX, SHA-256
% A48CB5CD1716974B686AC1CBA681CA4B17BC72F9043B78AD2528ACA41FCF814F.
% Exact logical German article SHA-256:
% A61962BD0D031352C9EB06FDD7A39813B49289B8361E22D5082F49B0AB0BD439.
% This file is not a Singapore, Taiwan, Hong Kong, or Macao localization.
\documentclass[11pt]{article}
\usepackage[a4paper,margin=2.35cm]{geometry}
\usepackage{fontspec}
\usepackage{xeCJK}
\usepackage{amsmath,amssymb}
\setmainfont{Noto Serif}
\setCJKmainfont[BoldFont=Microsoft YaHei]{SimSun}
\setCJKsansfont{Microsoft YaHei}
\XeTeXlinebreaklocale "zh"
\XeTeXlinebreakskip=0pt plus 1pt
\setlength{\parindent}{2em}
\setlength{\parskip}{0.45em}
\setlength{\emergencystretch}{3em}
\newcommand{\srcemph}[1]{\textbf{#1}}
\newcommand{\srctextsc}[1]{\textbf{#1}}
\begin{document}

\setcounter{footnote}{0}
\section*{33. 算术观点下的超复量与表示论}
\srcemph{Atti Congresso Bologna 2 (1928), S. 71--73.}

\begingroup
\renewcommand{\thefootnote}{(\arabic{footnote})}

\begin{center}
\textsc{Emmy Noether} (Göttingen -- Germania)\\[-0.1em]
\rule{3em}{0.4pt}\\[0.8em]
\srctextsc{超复量与表示论}\\
\srctextsc{（算术观点下）}\footnote{详细论述将发表于 Math. Zeitschr.；在那里的一篇续文还将讨论非交换伽罗瓦理论。关于“分裂域”，亦参见 R. Brauer 与 E. Noether, Ber. Berl. Ak. 1927；关于一般结构定理，参见 G. Köthe 在本卷 Atti 中的报告。}
\end{center}

我想说明，怎样把表示论——尤其是群表示论——理解为模的算子同构类与理想的算子同构类的理论，以及怎样把后两者纳入一个推广的群概念，即带算子的群。\footnote{带算子的群可追溯到 W. Krull (Math. Zeitschr. 23) 和 O. Schmidt (Math. Zeitschr. 29)。本文对给定表示之约化问题的处理，在表示除环为交换的限制下，也来自 W. Krull, Theorie und Anwendung der verallgemeinerten Abelschen Gruppen, Heidelberger Ber. 1926。}

所谓\srcemph{一个群的表示}——在预先给定的除环 $P$ 中——众所周知，是指如下情形：对元素 $a,b,\ldots$（它们属于群 $\mathfrak{G}$）分别配以矩阵 $A,B,\ldots$，矩阵的系数取自 $P$，使乘积 $ab$ 对应于乘积 $AB$；于是这个矩阵系统成为该群的一个同态像。与此同时，也给出了“群环”的一个表示；群环由所有“群数”（即群环元素）组成，也就是所有线性组合 $a\alpha+b\beta+\cdots+c\gamma$，其中 $a,b,\ldots,c$ 每次取有限多个群元素，而 $\alpha,\beta,\ldots,\gamma$ 相应地取 $P$ 中有限多个元素；这里把 $a,b,\ldots,c$ 视为线性无关，乘法由群乘法和运算法则定义。群环的表示同时关于加法同态；因此，群的表示被纳入\srcemph{环的表示}这一一般问题之下，在其中要求同时关于加法和乘法同态。

这里有两个本质问题。第一个是\srcemph{给定表示的约化}问题，以及在约化中出现的不可约组成部分是否唯一的问题；第二个则是询问一个给定环在一个给定的、不必交换的除环中的全部可能表示，例如只考虑其中的不可约表示。

第一个问题远为简单；这一点已经表现在，对于待表示的环以及通常非交换的表示除环，都不需要任何特殊假设。表示由固定阶数的矩阵给出这一事实，已经提供了一切必要的有限性假设。借助带算子的群的理论，这个问题可以完全处理；因此我先简短说明这一理论。

称群 $\mathfrak{B}$ 为一个\srcemph{带算子系的群}，如果同时给定一个符号系统 $\Theta,H,\ldots$——即算子系——并且表达式 $\Theta(a),H(a),\ldots$ 对每个 $a$（取自 $\mathfrak{B}$）都唯一确定，且给出 $\mathfrak{B}$ 中的一个元素，而且满足分配律 $\Theta(ab)=\Theta(a)\Theta(b)$。两个带有同一算子系的群称为算子同构，如果它们在通常意义下同构，并且由 $a$ 与 $\bar a$ 的对应还推出 $\Theta(a)$ 与 $\Theta(\bar a)$、$H(a)$ 与 $H(\bar a),\ldots$ 的对应。若只把自身也容许该算子系的子群算作可容许子群（因而所谓单群，是指除单位元外没有可容许的真正规子群），那么关于合成列的若尔当--赫尔德定理仍以如下形式成立：如果一个带算子系的群确有合成列，那么其因子群（合成因子）的系统在算子同构意义下唯一确定，至多差一个排列。在同一假设下，分解为直接不可分解因子的定理也成立，并且这种分解在算子同构意义下唯一。

它同表示问题的联系在于，理想和模尤其属于带算子的群：它们相对于加法被看作阿贝尔群，而算子由同环元素相乘给出。每一个表示都由一个表示模产生，也就是由一个关于该环和表示除环的模产生。更准确地说，每一个模类，即彼此算子同构的表示模所成的类，产生同一个表示。这样，表示约化的唯一性问题便由合成列定理和直积定理回答。把合成列应用于表示模，就对每个矩阵给出如下形式的约化：
\[
\begin{pmatrix}
B_1&0&0\\
&B_2&0\\
A_{ik}&\cdots&B_r
\end{pmatrix},
\]
其中对应于因子群的对角矩阵产生一个“不可约”表示，也就是说，它本身不再容许这样的约化。

既然这些因子群的类唯一确定，对角表示也同样唯一，这里的唯一性是在表示等价的意义下。类似地，直积定理给出“分解”为“不可分解”组成部分的唯一性：
\[
\begin{pmatrix}
C_1&&&\\
&C_2&&\\
&&\cdots&\\
&&&C_s
\end{pmatrix},
\]
其中每个 $C_i$ 随后还可以按合成列定理唯一地继续约化。

我想在超复系统的表示情形下，更一般地在满足“双链条件”的环的情形下，勾画\srcemph{第二个问题}。其结果是，每一个不可约（单）模的算子同构类都是某个理想的算子同构类；因而，由单侧理想组成的合成列已经通过其合成因子给出全部不可约表示。更准确地说，只须考察模去根基（在此即 Jacobson 根）所得商环中的合成列，而该商环成为一个“完全可约”环（现代称半单环）。问题由此化归为完全可约环的结构研究；因而，把自同态除环中的表示作为出发点就显得与问题相适应。这里所说的是如下事实：这样一个环的一个双侧单分量中的所有单右理想都属于同一个类，因此有同一个自同态环；由于这些理想是单的，这个自同态环成为除环，同时也是单左理想的自同态环。这个单环本身同自同态除环上的全体矩阵系统同构；而从这一表示出发，便产生关于一个子除环的一切表示，特别是关于该超复系统的系数域的一切表示，只要把自同态除环的元素本身替换为相应的矩阵。于是，在超复系统情形下已知的 Wedderburn 结构定理，在这里通过来自一般群论的自同态除环概念得到新的解释，并且同时被识别为超复系统表示论的来源。由此产生非交换除环的伽罗瓦理论中的新问题，它们同“分裂域”问题紧密相连。同时，通过不可分解理想的自同态环，也展望到一般的、并非完全可约的环的结构定理。群论，在它扩展到最一般的带算子的群之后，在这里也再次显现为一种整理原则。

\endgroup

\end{document}
